\documentclass[12pt]{article}
\usepackage[T1]{fontenc}
\usepackage[utf8]{inputenc}
\usepackage{lmodern}
\usepackage{textcomp}
\usepackage{enumitem}
\usepackage{geometry}
\geometry{margin=1in}
\begin{document}
\begin{center}
Answer Key - Political Science - Undergraduate Level Assessment 16 \
\small (Answers are ordered exactly as in the question paper.)
\end{center}

\section*{Multiple Choice}
\begin{enumerate}[label=\arabic*.]
\item \textbf{Answer:} D
\\ \textit{Options:} A) A constitutional amendment to the contrary of a court decision; B) The impeachment of a federal judge; C) A governor's failure to enforce a court decision; D) A national election recalling an unpopular judge
\\ \textit{Note:} A national election recalling an unpopular judge cannot limit the influence of the federal courts because federal judges are appointed for life and can only be removed through impeachment, not by popular vote or recall elections.
\item \textbf{Answer:} B
\\ \textit{Options:} A) the Senate.; B) the president.; C) the Secretary of State.; D) the chairman of the Joint Chiefs.
\\ \textit{Note:} The president has the constitutional authority to conduct foreign relations and formally recognize other nations, a power that is not granted to Congress.
\item \textbf{Answer:} A
\\ \textit{Options:} A) This would be an 'empire of liberty'; B) This empire would be multi-cultural; C) This type of empire would be based on expansion; D) This would be empire free of slavery
\\ \textit{Note:} Americans believed they could establish an 'empire of liberty' after 1776 because they envisioned a new type of governance that prioritized democratic ideals and individual freedoms, contrasting sharply with the imperial models based on conquest and oppression seen in European empires.
\item \textbf{Answer:} A
\\ \textit{Options:} A) Realism; B) Idealism; C) Liberalism; D) None of the above
\\ \textit{Note:} Realism posits that states are the principal entities in international relations, prioritizing their national interests and power in an anarchic global system.
\item \textbf{Answer:} A
\\ \textit{Options:} A) the Security Council.; B) the Chamber of Deputies.; C) the Council of Ministers.; D) the Secretariat.
\\ \textit{Note:} The Security Council holds real power within the United Nations because it is responsible for maintaining international peace and security, has the authority to make binding decisions, and can impose sanctions or authorize military action, which are not within the purview of other UN bodies.
\end{enumerate}

\section*{True / False}
\begin{enumerate}[label=\arabic*.]
\setcounter{enumi}{5}
\item \textbf{Answer:} False
\\ \textit{Options:} A) True; B) False
\\ \textit{Note:} The line-item veto was found unconstitutional because it violated the Presentment Clause of the Constitution by allowing the President to unilaterally amend or repeal parts of legislation, rather than delegating powers to the states.
\item \textbf{Answer:} False
\\ \textit{Options:} A) True; B) False
\\ \textit{Note:} Failed felony prosecutions at the state level cannot be retried federally on appeal because of the principle of double jeopardy, which prevents an individual from being tried for the same offense in multiple jurisdictions after an acquittal or conviction.
\item \textbf{Answer:} True
\\ \textit{Options:} A) True; B) False
\\ \textit{Note:} Entitlement programs, such as Social Security, Medicare, and Medicaid, are considered mandatory spending because they are legally required expenditures that automatically allocate funds to eligible individuals, making them the largest portion of the federal budget.
\item \textbf{Answer:} True
\\ \textit{Options:} A) True; B) False
\\ \textit{Note:} Critical Information Infrastructures are deemed the referent object in contemporary cyber-security because they are essential for national security, economic stability, and public safety, making their protection a priority in safeguarding against cyber threats.
\item \textbf{Answer:} False
\\ \textit{Options:} A) True; B) False
\\ \textit{Note:} Research has shown that individuals with higher education levels are often less supportive of government-enforced affirmative action programs, as they may prioritize merit-based criteria over race-based policies.
\end{enumerate}

\section*{Long Form Response}
\begin{enumerate}[label=\arabic*.]
\setcounter{enumi}{10}
\item \textbf{Answer:} The state plays a crucial role in achieving human security through various avenues such as governance, law enforcement, social services, and international relations. Effective governance ensures the establishment of a legal framework that protects citizens' rights and freedoms, fostering an environment of stability and trust. Law enforcement agencies maintain order and protect individuals from violence and crime, which directly contributes to personal security. Additionally, the state provides social services, including healthcare and education, which enhance the overall quality of life and empower individuals. Finally, through diplomatic efforts and international cooperation, the state can address transnational threats, contributing to the collective security of its citizens. These roles collectively imply that a strong and responsive state is fundamental to societal well-being, as it creates conditions that allow individuals to thrive in a safe and supportive environment.
\\ \textit{Note:} The answer correctly identifies the multifaceted roles of the state in promoting human security-through governance, law enforcement, social services, and international relations-which collectively enhance societal well-being by ensuring stability, protecting rights, and maintaining order.
\item \textbf{Answer:} The realist perspective on the international system posits that the anarchic nature of international relations, characterized by the absence of a central authority, drives states to prioritize their security and power. Realists argue that this competition for power and survival is a primary cause of war, as states may engage in conflict to assert dominance or respond to perceived threats. For example, the World Wars exemplify how national interests and power struggles can escalate into large-scale conflicts. Additionally, the ongoing tensions in regions like the South China Sea illustrate how states may resort to military action to secure their interests in an anarchic environment. Thus, realists contend that the intrinsic nature of international politics, driven by power dynamics, inherently leads to conflict.
\\ \textit{Note:} correct because it accurately describes the realist perspective's emphasis on the anarchic structure of the international system, the prioritization of power and security by states, and how these factors contribute to the outbreak of wars, illustrated by historical examples like the World Wars.
\end{enumerate}

\vfill
\noindent\textit{End of Answer Key}
\end{document}